ESPEON é uma ferramenta que faz uso majoritariamente da \textit{API chrome.tabs} para extrair dados de navegação do discente. Basicamente, mediante concessão de acesso, ilustrada no \textit{Script} \ref{scripts/manifest.json}, o \textit{software} consegue acessar o objeto \textit{tab}, que é uma dada aba do navegador. Conforme o conceito de objeto, na Programação Orientada a Objetos (POO), trata-se de uma abstração de algum elemento real, com atributos e métodos \cite{kimObjectorientedDatabasesDefinition1990}. O ESPEON faz a leitura desse objeto e processa seus valores no \textit{backend} - quando solicitada emissão de relatório - para construir relatórios de aula, parametrizados por tempo e estado. O tempo se transcreve em métricas de atenção, enquanto o estado está mais associado às métricas de engajamento, que são medidas pelo estado de periféricos de áudio e vídeo do aluno e o estado da atributos da aba, todos valores booleanos.

\spaceH{}
\lstinputlisting[
    language=json,
    caption={\textit{Manifest} do ESPEON},
    firstline=1,
    lastline=15
]{scripts/manifest.json}
\label{scripts/manifest.json}

O processo de leitura das abas ocorre quando há o disparo do gatilho, que é a troca de abas. Para definir a aba desejada, ou seja, a aula, o \textit{Service Worker} da extensão busca uma Expressão Regular (\textit{regex}) na propriedade \text{title} do objeto \textit{tab}; a Figura \ref{fig:regex} ilustra a composição da \textit{regex}.


\begin{figure}[H]
    \centering
    \begin{minipage}{0.7\textwidth}
        \centering
        
\includegraphics[width=.9\textwidth]{espeon/regex.drawio.png}
        \caption{Composição da \textit{Regex}}
        \label{fig:regex}
    \end{minipage}\hfill
\end{figure}

Para garantir a conformidade com essa expressão, o docente faz um cadastro prévio de sua aula no Banco de Dados, por meio do formulário disposto na Figura \ref{fig:formulario-aula}, no qual será gerada uma \textit{string} (ver Figura \ref{fig:clipboard}) com essa máscara, para que o docente, ao acessar a plataforma de \textit{Webconference} de sua escolha, preencha o título da reunião com essa expressão.

\begin{figure}[H]
    \centering
    \begin{minipage}{0.7\textwidth}
        \centering
        \includegraphics[width=.9\textwidth]{espeon/formulario.aula.png}
        \caption{Formulário de Aula}
        \label{fig:formulario-aula}
    \end{minipage}\hfill
\end{figure}

\begin{figure}[H]
    \centering
    \begin{minipage}{0.7\textwidth}
        \centering
        \includegraphics[width=.9\textwidth]{espeon/clipboard.png}
        \caption{Expressão Reguar Gerada}
        \label{fig:clipboard}
    \end{minipage}\hfill
\end{figure}

Então, quando o aluno acessa essa reunião por meio da extensão e o \textit{Service Worker} validar o título, ele estará sendo monitorado, sem qualquer necessidade de interação adicional. Suas interações com a guia e o estado dos periféricos serão armazenados em um banco de \textit{logs} para posterior composição dos relatórios.

% colocar figura com as permissões do manifest

%propriedades

%falar do armazenamento

Munido desses \textit{logs}, o \textit{backend} é capaz de processá-los em métricas de estatística descritiva. As métricas são devolvidas ao cliente por API nas formas de \textit{Javascript Object Notation (JSON)} e PDF. A primeira opção é particularmente versátil porque permite que docentes consumam as métricas em outras aplicações ou até mesmo em \textit{Software} Proprietário como o Power BI, para melhor visualizar os dados. Já as \textit{responses} em PDF foram projetadas como uma solução mais ágil de extração das métricas, bastando ao docente apenas ler o arquivo e interpretar os dados.